\subsection{Basic Exponential Models} 
In a \emph{linear} model, a constant change in one variable leads to a constant change in another. Linear functions show a \makeblank[1in] rate of change.
\begin{example}
You have \$20, and you get \$5 every day. We can write the amount of money you have, $m$, as a function of the number of days that have passed, $t$:
\[
m(t)=20+5t
\]
\end{example}
In an \emph{exponential} model, a constant change in one variable leads to a constant multiplication in another. In exponential models, the rate of change is \makeblank[1in] when the value of the function is \makeblank[1in], and \makeblank[1in] when the value of the function is \makeblank[1in].
\begin{example}
You have \$5, and your money doubles every day. We can write the amount of money you have, $m$, as a function of the number of days that have passed, $t$:
\[
m(t)=5\cdot 2^t
\]
\end{example}
\begin{displaybox}[Exponential Model (Basic Form)]
An exponential model is a function of the form
\[
f(x)=f_0\cdot m^x
\]
where $f_0$ is the starting value (the value of $f(x)$ when $x=0$) and $m$ is the multiplier (what $f(x)$ is multiplied by each time $x$ increases by 1).
\end{displaybox}
\begin{example}
A new web site is growing quickly in popularity: every week, the number of visitors triples. If the site has 10 visitors during is first week, write a function $v(t)$ showing the number of visitors $t$ weeks after the web site launches.
\begin{itemize}
\item The word ``triple'' means multiply by \makeblanksmall, so \makeblank[1in]
\item The starting value is \makeblanksmall (``the first week'' means \makeblanksmall full weeks have passed) so \makeblank[1in]
\end{itemize}
\[
v(t)=\makeblank
\]
\end{example}